\documentclass[11pt,a4paper]{article}

\usepackage[utf8]{inputenc}
\usepackage[T1]{fontenc}
\usepackage[english]{babel}
\usepackage[margin=2.5cm]{geometry}
\usepackage{amsmath}
\usepackage{amssymb}
\usepackage{amsthm}
\usepackage{mathtools}
\usepackage{microtype}
\usepackage{hyperref}
\usepackage{cleveref}

\hypersetup{
    colorlinks=true,
    linkcolor=blue,
    filecolor=magenta,      
    urlcolor=cyan,
    citecolor=red,
}

\title{\textbf{Analysis of $K^{(3)}$ Dependence on Network Depth $L$}}
\author{}
\date{}

\begin{document}
\maketitle


\tableofcontents

\newpage

\section{Framework and Notations}

In this section, we introduce the mathematical framework and notations used throughout this document, building upon the definitions established in the work on the Neural Tangent Hierarchy (NTH) \cite{dyer2019asymptoticswidenetworksfeynman}.

\subsection{Neural Network Architecture}

We consider a fully-connected feedforward neural network with $H$ hidden layers of width $m$.
\begin{itemize}
    \item The network input is a vector $x \in \mathbb{R}^d$.
    \item The output of layer $\ell$ is denoted by $x^{(\ell)}$. For $\ell \in \{1, \dots, H\}$, it is computed recursively as:
    \begin{equation}
        x^{(\ell)} = \frac{1}{\sqrt{m}}\sigma(W^{(\ell)}x^{(\ell-1)})
    \end{equation}
    where $x^{(0)} = x$ is the input, $W^{(\ell)}$ is the weight matrix of layer $\ell$, and $\sigma$ is a non-linear activation function applied component-wise.
    \item The final output of the network is a linear combination of the last hidden layer's output:
    \begin{equation}
        f(x, \theta) = a^T x^{(H)}
    \end{equation}
    where $a \in \mathbb{R}^m$ is the weight vector of the output layer.
    \item The set of all trainable parameters of the network is denoted $\theta = (\text{vec}(W^{(1)}), \dots, \text{vec}(W^{(H)}), a)$.
    \item The derivative of the activation function for layer $\ell$ is represented by a diagonal matrix $\sigma'_\ell(x) = \text{diag}(\sigma'(W^{(\ell)}x^{(\ell-1)}))$.
\end{itemize}

\subsection{The Neural Tangent Kernel (NTK)}

The network is trained via gradient descent on the quadratic loss: $L(\theta) = \frac{1}{2n}\sum_{\alpha=1}^n (f(x_\alpha, \theta) - y_\alpha)^2$. The dynamics of the network function $f(x, \theta_t)$ during continuous-time training (gradient flow) are governed by the Neural Tangent Kernel (NTK), denoted $K^{(2)}_t$.
\begin{equation}
\frac{d}{dt} f(x, \theta_t) = -\frac{1}{n}\sum_{\beta=1}^n K^{(2)}_t(x, x_\beta) (f(x_\beta, \theta_t) - y_\beta)
\end{equation}
The NTK is defined as the inner product of the gradients of the output function with respect to the network parameters:
\begin{equation}
K^{(2)}_t(x_\alpha, x_\beta) := \langle \nabla_\theta f(x_\alpha, \theta_t), \nabla_\theta f(x_\beta, \theta_t) \rangle
\end{equation}
It can be decomposed as a sum over the network layers:
\begin{equation}
K^{(2)}_t(x_\alpha, x_\beta) = \langle x^{(H)}_\alpha, x^{(H)}_\beta \rangle + \sum_{\ell=1}^{H} \langle G^{(\ell)}_\alpha, G^{(\ell)}_\beta \rangle \langle x^{(\ell-1)}_\alpha, x^{(\ell-1)}_\beta \rangle
\end{equation}
where the vector $G^{(\ell)}_\mu$ represents the backpropagated gradient up to layer $\ell$:
\begin{equation}
G^{(\ell)}_\mu = \frac{1}{\sqrt{m}}\sigma'_\ell(x_\mu) (W^{(\ell+1)})^T \cdots \frac{1}{\sqrt{m}}\sigma'_H(x_\mu) a_t
\end{equation}

\subsection{The Neural Tangent Hierarchy (NTH)}
For finite-width networks ($m < \infty$), the NTK $K^{(2)}_t$ is not constant but evolves during training. Its dynamics are dictated by the higher-order kernel $K^{(3)}_t$. This gives rise to an infinite sequence of coupled ordinary differential equations, known as the Neural Tangent Hierarchy (NTH):
\begin{equation}
\frac{d}{dt} K^{(r)}_t(x_{\alpha_1}, \dots, x_{\alpha_r}) = -\frac{1}{n} \sum_{\gamma=1}^n K^{(r+1)}_t(x_{\alpha_1}, \dots, x_{\alpha_r}, x_\gamma) (f_t(x_\gamma) - y_\gamma)
\end{equation}
valid for all $r \ge 2$. The higher-order kernel $K^{(r+1)}$ is defined recursively as the derivative of $K^{(r)}$ with respect to the parameters, contracted with the gradient of the network output:
\begin{equation}
K^{(r+1)}_t(x_1, \dots, x_{r+1}) := \langle \nabla_\theta K^{(r)}_t(x_1, \dots, x_r), \nabla_\theta f_t(x_{r+1}) \rangle
\end{equation}
This hierarchy captures the complete training dynamics beyond the infinite-width approximation where only $K^{(2)}$ is considered and is constant. The purpose of this document is to derive and analyze the explicit form of $K^{(3)}$, which represents the first correction to the NTK dynamics.

\section{Introduction}

This document analyzes the behavior of the entries of the $K^{(3)}$ tensor as a function of the network depth, denoted by $L$ (or $H$ in the code), for a fixed network width $m$. The analysis is based on the formulas from the Neural Tangent Hierarchy (NTH) framework. We aim to understand if the tensor entries decay as the depth increases, and specifically, if this decay is exponential.

\section{Scaling of Core Components in the NTK Formalism}

We consider a Multi-Layer Perceptron (MLP) of depth $L$ and uniform width $m$. The weights $W^{(\ell)}$ for $\ell=1, \dots, L$ are $m \times m$ matrices, and the output weights $a$ are in $\mathbb{R}^m$. We use the NTK scaling, where weights are initialized i.i.d. from a distribution with variance $c_W^2/m$ (e.g., $\mathcal{N}(0, c_W^2/m)$). The output layer is scaled by $1/\sqrt{m}$.

The forward activations $x^{(p)}_\mu$ and backward propagated vectors $G^{(\ell)}_\mu$ are the building blocks of the NTK.
The activations are defined as $x^{(p)}_\mu = \frac{1}{\sqrt{m}}\sigma(W^{(p)} x^{(p-1)}_\mu)$. With proper initialization ($c_W^2=2$ for ReLU), the norm of the activations is stable across layers:
\begin{equation}
    \mathbb{E}[\|x^{(p)}_\mu\|^2] \approx \|x^{(p-1)}_\mu\|^2
\end{equation}
Thus, we assume $\|x^{(p)}_\mu\| \sim O(1)$ for all layers $p$.

The backward vectors $G^{(\ell)}_\mu$ are defined recursively:
\begin{align}
    G^{(L)}_\mu &= \frac{a}{\sqrt{m}} \\
    G^{(\ell)}_\mu &= \frac{(W^{(\ell+1)})^T}{\sqrt{m}} \sigma'_{\ell+1}(x_\mu) G^{(\ell+1)}_\mu
\end{align}
where $\sigma'_{\ell+1}$ is the diagonal matrix of activation derivatives. A crucial point is to analyze the norm of the propagation operator $\mathcal{P}_\ell = \frac{(W^{(\ell+1)})^T}{\sqrt{m}} \sigma'_{\ell+1}(x_\mu)$. Based on random matrix theory, for a large matrix $W$ with i.i.d. entries of variance $\sigma_w^2$, its operator norm $\|W\|_{\text{op}}$ converges to $2\sigma_w\sqrt{m}$. Thus, the norm of our scaled operator is:
\begin{equation}
    \|\mathcal{P}_\ell\|_{\text{op}} \le \left\| \frac{W^{(\ell+1)}}{\sqrt{m}} \right\|_{\text{op}} \xrightarrow{m\to\infty} 2\sigma_w
\end{equation}
For standard initialization schemes (e.g., He initialization with $\sigma_w^2=2/m_{in}$), this value is $O(1)$ but not necessarily less than 1. A naive application of the sub-multiplicative property might suggest an exponential explosion of norms, as $\|G^{(\ell)}\| \le (2\sigma_w)^{L-\ell} \|G^{(L)}\|$.

However, this reasoning is flawed because we are multiplying a sequence of \textbf{independent} random matrices. The vector $G^{(\ell+1)}_\mu$, after being transformed by $\mathcal{P}_\ell$, is not aligned with the direction of maximal amplification of the next operator $\mathcal{P}_{\ell-1}$. The theory of products of random matrices shows that for deep networks, the norm is preserved on average if the initialization is chosen correctly to achieve a state of "dynamical isometry". This ensures that the system avoids both exponential explosion and vanishing of the signal. Therefore, we can confidently assume that the norms are stable across layers:
\begin{equation}
    \|G^{(\ell)}_\mu\| \sim O(1) \quad \text{for all } \ell
\end{equation}
These vectors do not exhibit a systematic decay or growth with depth.

\section{Analysis of the $K^{(3)}$ Tensor}

The $K^{(3)}$ tensor is constructed from the directional derivatives of $x^{(p)}_\mu$ and $G^{(\ell)}_\mu$, denoted $\delta_\gamma x^{(p)}_\mu$ and $\delta_\gamma G^{(\ell)}_\mu$.

\subsection{Forward Derivative Term $\delta_\gamma x^{(p)}_\mu$}

The recursive formula for $\delta_\gamma x^{(p)}_\mu$ is:
\begin{equation}
    \delta_\gamma x^{(p)}_\mu = \frac{1}{\sqrt{m}} \sigma'_{p}(x_\mu) \left( W^{(p)} (\delta_\gamma x^{(p-1)}_\mu) + \langle x^{(p-1)}_\gamma, x^{(p-1)}_\mu \rangle G^{(p)}_\gamma \right)
\end{equation}
with $\delta_\gamma x^{(0)}_\mu = 0$. Analyzing the norm:
\begin{equation}
    \|\delta_\gamma x^{(p)}_\mu\| \le \frac{1}{\sqrt{m}} \|\sigma'_{p}\| \left( \|W^{(p)}\| \|\delta_\gamma x^{(p-1)}_\mu\| + |\langle \dots \rangle| \|G^{(p)}_\gamma\| \right)
\end{equation}
Using $\|W^{(p)}\| \sim O(\sqrt{m})$ and that other norms are $O(1)$, we get:
\begin{equation}
    \|\delta_\gamma x^{(p)}_\mu\| \lesssim \|\delta_\gamma x^{(p-1)}_\mu\| + O(1/\sqrt{m})
\end{equation}
Starting from $\|\delta_\gamma x^{(0)}_\mu\| = 0$, we find $\|\delta_\gamma x^{(p)}_\mu\| \sim O(p/\sqrt{m})$. This indicates a linear growth with layer index $p$, scaled by $1/\sqrt{m}$.

\subsection{Backward Derivative Term $\delta_\gamma G^{(\ell)}_\mu$}

The term $\delta_\gamma G^{(\ell)}_\mu$ has a more complex structure, involving a sum over subsequent layers $p$ from $\ell$ to $L$. A representative term in this sum is (ignoring propagators for simplicity):
\begin{equation}
    \text{term}_p \sim \frac{1}{m} G^{(p+1)}_\gamma \langle x^{(p)}_\gamma, x^{(p)}_\mu \rangle
\end{equation}
This term has a factor of $1/m$. The propagation of this term back to layer $\ell$ involves products of matrices of the form $(W^{(k)})^T/\sqrt{m}$, which have $O(1)$ norm. Therefore, each term in the sum for $\delta_\gamma G^{(\ell)}_\mu$ is of order $O(1/m)$. Since there are $L-\ell$ such terms, we get:
\begin{equation}
    \|\delta_\gamma G^{(\ell)}_\mu\| \sim O\left(\frac{L-\ell}{m}\right)
\end{equation}
This term shows a linear dependency on the number of subsequent layers and is scaled by $1/m$.

\subsection{Overall Scaling of $K^{(3)}$}

The full expression for $K^{(3)}_{\alpha\beta\gamma}$ is a sum of terms involving scalar products of the quantities analyzed above. For example:
\begin{equation}
K^{(3)}_{\alpha\beta\gamma} = \langle \delta_\gamma x^{(L)}_\alpha, x^{(L)}_\beta \rangle + \langle x^{(L)}_\alpha, \delta_\gamma x^{(L)}_\beta \rangle + \dots
\end{equation}
The dominant term comes from the forward derivative part:
\begin{equation}
    \langle \delta_\gamma x^{(L)}_\alpha, x^{(L)}_\beta \rangle \sim O(L/\sqrt{m})
\end{equation}
The other terms involving $\delta_\gamma G$ are smaller, of order $O(L/m)$. Thus, for large $m$:
\begin{equation}
    K^{(3)}_{\alpha\beta\gamma} \sim O(L/\sqrt{m})
\end{equation}

\section{Conclusion on Depth Dependence}

Our analysis shows that the magnitude of the entries of the $K^{(3)}$ tensor does \textbf{not} decay with depth $L$. Instead, it appears to grow linearly with $L$. The scaling factor is $1/\sqrt{m}$.

This contradicts the initial suspicion of an exponential decay. The key observations are:
\begin{itemize}
    \item The fundamental vectors $x^{(p)}$ and $G^{(\ell)}$ have norms that are stable across layers and do not decay.
    \item The derivative terms $\delta_\gamma x^{(p)}$ and $\delta_\gamma G^{(\ell)}$ introduce factors of $1/\sqrt{m}$ and $1/m$ respectively, but their recursive/summative nature leads to a linear growth with the number of layers involved ($p$ or $L-\ell$).
\end{itemize}

Therefore, for a fixed width $m$, we expect the values of $K^{(3)}$ to increase with the depth $L$. There is no evidence of an exponential decay. The factor $1/\sqrt{m}$ controls the overall magnitude but does not introduce a decay with depth.


\newpage

\section{Detailed Scaling Analysis with Respect to Depth $H$}

We now provide a more detailed analysis of the scaling of $K^{(3)}_{\alpha\beta\gamma}$ with respect to the network depth $H$ (denoted $L$ in some contexts) and width $m$. This analysis is based on the fully expanded formula and relies on standard assumptions in the study of wide neural networks.

\subsection{Core Assumptions and Scaling of Components}
Our analysis rests on the following scaling assumptions, which are direct consequences of the NTK parameterization and standard random matrix theory results:
\begin{itemize}
    \item \textbf{Activations and G-vectors:} The norms of forward activations and backward-propagated vectors are stable across layers: $\|x^{(p)}_\mu\| \sim O(1)$ and $\|G^{(\ell)}_\mu\| \sim O(1)$.
    \item \textbf{Propagators:} Due to the principle of dynamical isometry targeted by modern initialization schemes, the operator norms of the forward and backward propagators are of order one, regardless of the number of matrices in the product: $\|\mathcal{J}^{(p \to j)}_\mu\|_{\text{op}} \sim O(1)$ and $\|\mathcal{B}^{(\ell \to p)}_\mu\|_{\text{op}} \sim O(1)$.
    \item \textbf{Source Terms:} The scaling of the source terms can be readily determined:
    \begin{itemize}
        \item Forward source term: $\mathcal{T}^{(j)}_{\gamma\mu} = \frac{\sigma'_{j}(x_\mu)}{\sqrt{m}} \langle x^{(j-1)}_\gamma, x^{(j-1)}_\mu \rangle G^{(j)}_\gamma$. Its norm scales as $\|\mathcal{T}^{(j)}_{\gamma\mu}\| \sim O(1/\sqrt{m})$.
        \item Backward source term: $\mathcal{U}^{(p)}_{\gamma\mu} = \frac{1}{m} G^{(p+1)}_\gamma \langle x^{(p)}_\gamma, x^{(p)}_\mu \rangle$. Its norm scales as $\|\mathcal{U}^{(p)}_{\gamma\mu}\| \sim O(1/m)$.
    \end{itemize}
\end{itemize}

\subsection{Term-by-Term Analysis}
We analyze the four main components of the fully expanded formula for $K^{(3)}_{\alpha\beta\gamma}$. For simplicity, we analyze one of the two symmetric terms in each summation.

\textbf{1. Output Term ($K^{(3,\text{out})}$)}
This term is given by the sum $\sum_{j=1}^{H} \langle \mathcal{J}^{(H \to j)}_{\alpha} \mathcal{T}^{(j)}_{\gamma\alpha}, x^{(H)}_\beta \rangle$.
The magnitude of each summand is:
\[
|\langle \mathcal{J}^{(H \to j)}_{\alpha} \mathcal{T}^{(j)}_{\gamma\alpha}, x^{(H)}_\beta \rangle| \le \|\mathcal{J}^{(H \to j)}_{\alpha}\|_{\text{op}} \cdot \|\mathcal{T}^{(j)}_{\gamma\alpha}\| \cdot \|x^{(H)}_\beta\| \sim O(1) \cdot O(1/\sqrt{m}) \cdot O(1) = O(1/\sqrt{m})
\]
Since we are summing $H$ such terms, the total contribution of this component is:
\[
K^{(3,\text{out})} \sim H \cdot O(1/\sqrt{m}) = O(H/\sqrt{m})
\]

\textbf{2. Backward Term from Weights ($K^{(3,G,W)}$)}
This term is the double summation $\sum_{\ell=1}^{H} \langle x^{(\ell-1)}_\alpha, x^{(\ell-1)}_\beta \rangle \sum_{p=\ell}^{H-1} \langle \mathcal{B}^{(\ell \to p)}_{\alpha} \mathcal{U}^{(p)}_{\gamma\alpha}, G^{(\ell)}_\beta \rangle$.
The magnitude of each summand in the inner loop is:
\[
|\langle \dots \rangle \langle \dots \rangle| \le |\langle x^{(\ell-1)}_\alpha, x^{(\ell-1)}_\beta \rangle| \cdot \|\mathcal{B}^{(\ell \to p)}_{\alpha}\|_{\text{op}} \cdot \|\mathcal{U}^{(p)}_{\gamma\alpha}\| \cdot \|G^{(\ell)}_\beta\| \sim O(1) \cdot O(1) \cdot O(1/m) \cdot O(1) = O(1/m)
\]
This double sum contains approximately $\sum_{\ell=1}^H (H-\ell) \approx H^2/2$ terms. Thus, the total contribution is:
\[
K^{(3,G,W)} \sim H^2 \cdot O(1/m) = O(H^2/m)
\]

\textbf{3. Backward Term from Output Layer ($K^{(3,G,a)}$)}
This term corresponds to the derivative with respect to the output weights $a_t$: $\sum_{\ell=1}^{H} \langle x^{(\ell-1)}_\alpha, x^{(\ell-1)}_\beta \rangle \langle \mathcal{B}^{(\ell \to H-1)}_{\alpha} \frac{x^{(H)}_\gamma}{\sqrt{m}}, G^{(\ell)}_\beta \rangle$.
The magnitude of each summand is:
\[
|\langle \dots \rangle \langle \dots \rangle| \le |\langle x^{(\ell-1)}_\alpha, x^{(\ell-1)}_\beta \rangle| \cdot \|\mathcal{B}^{(\ell \to H-1)}_{\alpha}\|_{\text{op}} \cdot \frac{\|x^{(H)}_\gamma\|}{\sqrt{m}} \cdot \|G^{(\ell)}_\beta\| \sim O(1) \cdot O(1) \cdot O(1/\sqrt{m}) \cdot O(1) = O(1/\sqrt{m})
\]
Summing $H$ such terms, the total contribution is:
\[
K^{(3,G,a)} \sim H \cdot O(1/\sqrt{m}) = O(H/\sqrt{m})
\]

\textbf{4. Forward Term from Activations ($K^{(3,x)}$)}
This is the final double summation: $\sum_{\ell=1}^{H} \langle G^{(\ell)}_\alpha, G^{(\ell)}_\beta \rangle \sum_{j=1}^{\ell-1} \langle \mathcal{J}^{(\ell-1 \to j)}_{\alpha} \mathcal{T}^{(j)}_{\gamma\alpha}, x^{(\ell-1)}_\beta \rangle$.
The magnitude of each summand in the inner loop is:
\[
|\langle \dots \rangle \langle \dots \rangle| \le |\langle G^{(\ell)}_\alpha, G^{(\ell)}_\beta \rangle| \cdot \|\mathcal{J}^{(\ell-1 \to j)}_{\alpha}\|_{\text{op}} \cdot \|\mathcal{T}^{(j)}_{\gamma\alpha}\| \cdot \|x^{(\ell-1)}_\beta\| \sim O(1) \cdot O(1) \cdot O(1/\sqrt{m}) \cdot O(1) = O(1/\sqrt{m})
\]
This double sum contains approximately $\sum_{\ell=1}^H (\ell-1) \approx H^2/2$ terms. The total contribution is therefore:
\[
K^{(3,x)} \sim H^2 \cdot O(1/\sqrt{m}) = O(H^2/\sqrt{m})
\]

\subsection{Overall Scaling and Conclusion}
Combining the four components, we have:
\[
K^{(3)}_{\alpha\beta\gamma} = O(H/\sqrt{m}) + O(H^2/m) + O(H/\sqrt{m}) + O(H^2/\sqrt{m})
\]
For a sufficiently large width $m$, the term $O(H^2/\sqrt{m})$ will dominate the term $O(H^2/m)$. Therefore, the overall scaling behavior of the third-order kernel is dictated by the $K^{(3,x)}$ term:
\begin{equation}
    K^{(3)}_{\alpha\beta\gamma} \sim O\left(\frac{H^2}{\sqrt{m}}\right)
\end{equation}
This detailed analysis confirms that the magnitude of the $K^{(3)}$ tensor entries does not decay with depth. Instead, it exhibits a \textbf{quadratic growth} with the depth $H$. This implies that the dynamics of the NTK itself become more pronounced and complex in very deep networks, a key feature of the finite-width correction captured by the Neural Tangent Hierarchy.

\newpage

\section{Implementation of the NTK, and the code correspondence}

The Python script \texttt{kernel3\_empirical.py} in the \href{https://github.com/janisaiad/DeNN-NTK}{github repository} provides a direct, empirical implementation of the analytical formulas for the third-order kernel, $K^{(3)}$, as derived within the Neural Tangent Hierarchy framework. This implementation translates the mathematical recursions and summations into executable code.

\subsection{Code Structure and Correspondence to Formulas}
The core logic is encapsulated within the `Kernel3Empirical` class. Its methods directly correspond to the mathematical objects we have analyzed:

\begin{itemize}
    \item \textbf{\texttt{\_compute\_G(ell, mu)}}: This helper function computes the backward-propagated vector $G^{(\ell)}_\mu$. It implements the standard recursive definition:
    \[ G^{(\ell)}_\mu = \frac{(W^{(\ell+1)})^T}{\sqrt{m}} \sigma'_{\ell+1}(x_\mu) G^{(\ell+1)}_\mu \]
    The recursion starts from the base case $G^{(H)}_\mu = a_t / \sqrt{m}$.

    \item \textbf{\texttt{\_compute\_delta\_x(p, gamma, mu)}}: This function calculates the forward derivative term $\delta_\gamma x^{(p)}_\mu$. It follows the recursive formula derived from the chain rule for forward activations:
    \[ \delta_\gamma x^{(p)}_\mu = \frac{1}{\sqrt{m}} \sigma'_{p}(x_\mu) \left( W^{(p)} (\delta_\gamma x^{(p-1)}_\mu) + \langle x^{(p-1)}_\gamma, x^{(p-1)}_\mu \rangle G^{(p)}_\gamma \right) \]
    The recursion starts from the base case $\delta_\gamma x^{(0)}_\mu = 0$.

    \item \textbf{\texttt{\_compute\_delta\_G(ell, gamma, mu)}}: This method implements the more complex formula for the backward derivative term $\delta_\gamma G^{(\ell)}_\mu$. It constructs the term by summing the contributions from the derivatives of all subsequent weights ($W^{(p+1)}$ for $p \ge \ell$) and the output layer weights ($a_t$). Each contribution is then propagated back to layer $\ell$ through a series of matrix-vector products, as specified by the analytical formula.

    \item \textbf{\texttt{kernel3(alpha, beta, gamma)}}: This is the main public method. It computes the final scalar value of $K^{(3)}_{\alpha\beta\gamma}$ by assembling all the pieces. It mirrors the high-level decomposition:
    \[ K^{(3)}_{\alpha\beta\gamma} = K^{(3, \text{out})}_{\alpha\beta\gamma} + \sum_{\ell=1}^H \left( K^{(3, G)}_{\alpha\beta\gamma, \ell} + K^{(3, x)}_{\alpha\beta\gamma, \ell} \right) \]
    It calls the helper methods to compute the necessary $\delta_\gamma x$, $\delta_\gamma G$, and $G$ vectors and then combines them through the appropriate inner products.
\end{itemize}

\subsection{Future Work: Optimization with Automatic Differentiation}

The current implementation is a literal translation of the mathematical formulas. While valuable for its clarity and direct correspondence to the theory, this manual approach has drawbacks: it is computationally intensive due to explicit recursion and can be complex to debug and extend.

In the future, a more efficient and robust approach will be pursued by leveraging modern automatic differentiation (AD) frameworks, particularly JAX. The vision is to move away from implementing the analytical derivatives manually and instead let the framework compute them automatically. This strategy is at the heart of libraries like Google's \texttt{neural-tangents}.

The plan is as follows:
\begin{enumerate}
    \item Define the second-order kernel, $K^{(2)}(x_\alpha, x_\beta)$, as a pure JAX function that depends on the network parameters $\theta$.
    \item Use JAX's function transformations, such as \texttt{jax.grad}, to compute the gradient of the kernel with respect to the parameters, $\nabla_\theta K^{(2)}_{\alpha\beta}$.
    \item The third-order kernel, $K^{(3)}$, is defined as the inner product $K^{(3)}_{\alpha\beta\gamma} = \langle \nabla_\theta K^{(2)}_{\alpha\beta}, \nabla_\theta f_\gamma \rangle$. Since $\nabla_\theta f_\gamma$ can also be computed automatically with \texttt{jax.grad}, the entire $K^{(3)}$ tensor can be obtained without manually writing out its complex formula.
\end{enumerate}

This AD-based approach promises significant benefits:
\begin{itemize}
    \item \textbf{Efficiency:} JAX can compile the entire computation into a highly optimized computational graph for execution on accelerators like GPUs or TPUs.
    \item \textbf{Simplicity and Robustness:} The code becomes much simpler, as we only need to define the base objects ($f$ and $K^{(2)}$), not their derivatives. This drastically reduces the risk of implementation errors.
    \item \textbf{Extensibility:} This methodology would make it far easier to compute even higher-order kernels, such as $K^{(4)}$, which would be practically infeasible to implement manually.
\end{itemize}

\newpage

\section{Empirical Validation: Experimental Setup and Analysis}

The script \texttt{kernel3\_vs\_depth\_finite\_empirical\_largeexperiments.py} is designed to empirically validate the theoretical scaling laws for $K^{(3)}$ derived in the previous sections. The primary goal is to investigate how the magnitude of the $K^{(3)}$ tensor entries depends on the network depth $H$ and width $M$.

\subsection{Experimental Design}

The experiment is structured as a systematic grid search over a range of network configurations. The main loop iterates through different values for:
\begin{itemize}
    \item \texttt{N\_VALUES}: The number of data points.
    \item \texttt{D\_IN\_VALUES}: The input feature dimension.
    \item \texttt{M\_VALUES}: The network width.
    \item \texttt{L\_VALUES}: The network depth (referred to as $H$ in our formulas).
\end{itemize}

For each configuration, the script performs the following steps:
\begin{enumerate}
    \item \textbf{Initialization}: It generates random, normalized input data using \texttt{generate\_data} and initializes the network weights according to a standard normal distribution via \texttt{init\_network}.
    \item \textbf{Forward Pass}: It calls \texttt{compute\_features\_and\_derivatives} to perform a full forward pass, collecting all the necessary intermediate values: the feature maps ($x^{(p)}$) and the diagonal matrices of activation derivatives ($\sigma'_p$) for each layer.
    \item \textbf{Kernel Computation}: An instance of the \texttt{Kernel3Empirical} class is created with the network weights and the results from the forward pass. The script then computes the entire $N \times N \times N$ tensor for $K^{(3)}$ by calling the \texttt{kernel3(i, j, k)} method for all unique triplets of indices, leveraging the tensor's symmetry.
\end{enumerate}

\subsection{Data Analysis and Connection to Theory}

Once the empirical \texttt{k3\_tensor} is computed, the script extracts quantitative metrics to analyze its properties.

\begin{itemize}
    \item \textbf{Infinity Norm}: The script computes \texttt{inf\_norm = M * np.max(np.abs(k3\_tensor))}. The purpose of this analysis is to study the kernel's asymptotic behavior. By scaling the empirical result with a factor depending on the width $M$, we can check if the quantity converges to a finite, depth-dependent limit as $M \to \infty$. Our theoretical analysis predicted that the entries of $K^{(3)}$ scale as $O(H^2/\sqrt{M})$. Therefore, multiplying the maximum absolute value by $\sqrt{M}$ should result in a quantity that grows quadratically with depth $H$. The script's use of a scaling factor of $M$ may be intended to test a different theoretical hypothesis or scaling regime.

    \item \textbf{Eigenvalue Analysis}: The script also performs a spectral analysis by examining the eigenvalues of 2D slices of the tensor (\texttt{k3\_tensor[:, :, k]}). This provides deeper insight into the operator's properties than its maximum entry alone. The mean and standard deviation of these eigenvalues are computed across all slices to produce a statistical summary of the kernel's spectrum. These spectral properties are also scaled by the width $M$.
\end{itemize}

The data generated by this script (the computed norms and eigenvalue statistics for varying $H$ and $M$) is saved and can be plotted to visually confirm or challenge the $O(H^2/\sqrt{M})$ scaling law derived theoretically. This provides a crucial bridge between the analytical formulas of the NTH and their concrete implications for the behavior of finite-width neural networks.

\newpage

\section{Deep Dive into Quantitative Metrics}

To bridge the gap between our theoretical analysis and the empirical results from the script, it's essential to understand precisely what the computed metrics represent and how they connect to the underlying theory.

\subsection{Analysis of the Scaled Infinity Norm}

The script computes \texttt{inf\_norm = M * np.max(np.abs(k3\_tensor))}. Let's break this down.

\begin{itemize}
    \item \textbf{What is the Infinity Norm?} The infinity norm of the $K^{(3)}$ tensor, denoted $\|K^{(3)}\|_\infty$, is simply the largest absolute value among all its entries:
    \[ \|K^{(3)}\|_\infty = \max_{\alpha, \beta, \gamma} |K^{(3)}_{\alpha\beta\gamma}| \]
    This metric provides a simple way to quantify the overall magnitude of the kernel.

    \item \textbf{The Purpose of Scaling by Width $M$}: In the infinite-width limit ($M \to \infty$), the standard NTK theory shows that $K^{(3)}$ vanishes, which is why the kernel dynamics are ignored. Our analysis for finite-width networks aims to characterize precisely \textit{how fast} it vanishes. The scaling factor is introduced to cancel out the width-dependent decay, thereby isolating the depth-dependent behavior.

    \item \textbf{Testing Theoretical Predictions}: Our detailed scaling analysis predicted that the dominant term in $K^{(3)}$ scales as $O(H^2/\sqrt{M})$. To empirically test this hypothesis, one should analyze the quantity $\sqrt{M} \cdot \|K^{(3)}\|_\infty$. If the theory is correct, this scaled value should converge to a limit that grows quadratically with depth ($H^2$) as $M$ becomes large.

    The experiment, by computing $M \cdot \|K^{(3)}\|_\infty$, is implicitly testing an alternative hypothesis where the kernel entries scale as $O(1/M)$. If our $O(1/\sqrt{M})$ theory is correct, the quantity computed in the script should not converge but instead grow like $\sqrt{M}$. By observing which of these scaled quantities (or another) stabilizes as $M$ increases, the experiment can empirically determine the true scaling law, thus validating or falsifying different theoretical claims.
\end{itemize}

\subsection{Spectral Analysis of Tensor Slices}

The eigenvalue analysis provides a much deeper view into the structure of $K^{(3)}$ than the infinity norm alone. The dynamics of the NTK are given by the equation:
\[ \frac{d}{dt} K^{(2)}_t = -\frac{1}{n} \sum_{\gamma=1}^n (f_t(x_\gamma) - y_\gamma) \cdot K^{(3)}_{t, \gamma} \]
where $K^{(3)}_{t, \gamma}$ is the $N \times N$ matrix slice of the tensor for a fixed index $\gamma$. The evolution of $K^{(2)}$ is thus a linear combination of these slice matrices.

\begin{itemize}
    \item \textbf{Physical Meaning of Eigenvalues}: The eigenvalues of a slice matrix $K^{(3)}_\gamma$ characterize how an error at a single data point, $x_\gamma$, influences the geometry of the feature space (as captured by $K^{(2)}$). An eigenvector of this matrix represents a direction in the input space, and its corresponding eigenvalue represents the magnitude and sign of the change in that direction. A large positive eigenvalue means that the kernel's value for pairs of points aligned with that eigenvector will increase, effectively pushing those representations further apart.

    \item \textbf{Mean and Standard Deviation}:
    \begin{itemize}
        \item \texttt{mean\_eigenvalues}: By averaging the eigenvalues of all slices, we obtain the "average" spectral signature of the $K^{(3)}$ operator. This tells us if there is a systematic, data-independent drift in the NTK. For instance, a consistently negative smallest eigenvalue across all slices would imply that gradient descent tends to contract the feature space along a specific direction, regardless of which data point has a large error.
        \item \texttt{std\_eigenvalues}: The standard deviation measures how much this spectral impact varies depending on the data point $x_\gamma$ driving the update. A small standard deviation would imply that the kernel evolves in a very predictable, almost data-independent manner. Conversely, a large standard deviation suggests that the geometry of the kernel is being sculpted in a highly complex, data-dependent way, with each training example pulling the kernel in a different "direction".
    \end{itemize}

    \item \textbf{Asymptotic Behavior}: As with the infinity norm, the eigenvalues are scaled by the width $M$ to study their asymptotic behavior. The eigenvalues of a matrix with entries scaling as $\epsilon$ will also scale as $\epsilon$. Therefore, if the entries of $K^{(3)}$ scale as $O(1/\sqrt{M})$, its eigenvalues should as well. The experiment provides the necessary data to check if the scaled spectrum converges to a stable, non-trivial limit as the width grows.
\end{itemize}


\end{document}
